\section{Uniform correction of the real moment map}\label{sec:balancing}

This section proves the estimate that permits independent degeneration
scales. It applies to the full complex equation $\cQ$, including its
lower-rank points. Compatibility of the corrected form with the
coefficient complex structure is not assumed.

\begin{theorem}\label{thm:balancing}
Let $\Omega_t$ be a smooth family of closed $K$-invariant forms on
a ball in $\cV$, with centered moment maps $\rho_t$ in the convention
of~\eqref{eq:form-error}, and suppose
\begin{equation}\label{eq:abstract-form-bound}
\|\Omega_{t,x}-\omega^{\rm amb}\|\le C(|t|+|x|^2).
\end{equation}
For every sufficiently small $x=(A,B)\in\cQ$ with $\mu(x)=\tau I$,
there is a Hermitian matrix $S$ such that
\begin{equation}\label{eq:balanced-correction}
T_Sx=(Ae^S,e^{-S}B)=y,\qquad
\rho_t(y)=\tau I,\qquad
\|S\|\le C'(|t|+|x|^2).
\end{equation}
At $\tau=0$, $S$ may be taken zero on the common kernel of $A$ and
$B^*$, with no off-diagonal block to its orthogonal complement.

The map on compact orbits defined by the common complex orbit is unique,
continuous through rank change, and smooth on each fixed-rank piece.
It covers all sufficiently small points of the corrected real level,
after a fixed adjustment of the domain radius. No constant depends on
the smaller positive eigenvalue of $A^*A+BB^*$.
\end{theorem}

\begin{proof}
Write $E=\rho_t-\mu$. On a ball where
$\|\Omega_t-\omega^{\rm amb}\|\le\epsilon$, every orthogonal
projection $P$ on $\CC^2$ satisfies
\begin{equation}\label{eq:projection-error}
|\tr(PE(x))|\le
\epsilon\bigl(\|AP\|^2+\|PB\|^2\bigr).
\end{equation}
To see this, decompose $x=n+u$ into the fixed and moving parts for
the circle generated by $P$:
\[
n=(A(1-P),(1-P)B),\qquad u=(AP,PB).
\]
Each moment component is constant on the connected fixed linear
subspace and has value zero there. On $n+su$,
$\xi_P(n+su)=s\xi_P(u)$ and $\|\xi_P(u)\|=\|u\|$. Hence
\[
\tr(PE(x))
=2\int_0^1s(\Omega_{t,n+su}-\omega^{\rm amb})
                 (\xi_P(u),u)\,ds,
\]
which proves~\eqref{eq:projection-error}. The segment stays in the
same ball. Testing rank-one projections also gives the matrix inequality
$-\epsilon M\le E\le\epsilon M$, where $M=A^*A+BB^*$.
This relative error is stronger than an additive bound near rank loss.

Interpolate $\Omega_{t,s}=(1-s)\omega^{\rm amb}+s\Omega_t$ and
$\rho_{t,s}=\mu+sE$ for $0\le s\le1$. They tame $I$ on a small
common ball. Suppose first $\tau\ne0$, so the action at $x$ is free.
Starting from $S=0$ at $s=0$, solve
$\rho_{t,s}(T_Sx)=\tau I$. The derivative in imaginary orbit
directions is invertible, because for every nonzero Hermitian $h$,
\[
d\tr(\rho_{t,s}h)(I\xi_h)
=2\Omega_{t,s}(\xi_h,I\xi_h)>0.
\]
This positive quadratic pairing gives injectivity; equal dimensions give
invertibility, without requiring symmetry of the derivative.
Positive polar coordinates give the same invertibility. More explicitly,
the Hermitian part of the logarithmic derivative of $e^S$ has, in an
eigenbasis of $S$, coefficients
$\sinh(\lambda_i-\lambda_j)/(\lambda_i-\lambda_j)$, interpreted as
one on the diagonal. They are positive. The remaining compact part
has zero moment derivative at a central value. This argument asserts
invertibility at each free point, not a uniform inverse bound.

The uniform bound instead follows from~\eqref{eq:projection-error}.
Let $P$ be the spectral projection onto a unit eigenvector of $S$,
with eigenvalue $\lambda$, and put $a=\|AP\|^2$, $b=\|PB\|^2$
at the original balanced point. Thus $a-b=\tau$. At a solution,
\begin{equation}\label{eq:eigen-balance}
0=(e^{2\lambda}-1)a+(1-e^{-2\lambda})b
        +s\tr(PE(T_Sx)).
\end{equation}
The last term is at most
$\epsilon(e^{2\lambda}a+e^{-2\lambda}b)$ in absolute value.
For $0<\epsilon<1/10$, both lower-bound coefficients are positive
when $\lambda>2\epsilon$, and both upper-bound coefficients are
negative when $\lambda<-2\epsilon$. Since $a+b>0$, neither is
possible. Thus $\|S\|\le2\epsilon$.

Use the ball of radius $2|x|$, with
$\epsilon=C(|t|+4|x|^2)$. For small $t,x$, the bound ensures
$|T_Sx|\le e^{2\epsilon}|x|<2|x|$. A continuation branch cannot
leave the domain or escape in $S$. Compactness gives a solution at
each limiting interpolation time, and the invertible linearization
continues it. This proves~\eqref{eq:balanced-correction}.

If $\tau=0$, flat balance gives $A^*A=BB^*=P_0$. Restrict to its
support $V_0$. Both maps vanish on the complementary kernel, and
the $U(V_0)$ action on the support is free. Apply the same argument
there. Equivariance under $U(V_0^\perp)$ makes the moment block
diagonal; its kernel block is zero by~\eqref{eq:projection-error},
and the off-diagonal block vanishes by invariance. The restricted
solution therefore solves the full equation. At $x=0$, take $S=0$.

To prove uniqueness of the compact orbit, take two small corrected
balanced points in the same complex orbit. After a polar decomposition
and compact change they are $z,T_Hz$ for Hermitian $H$. The squared
norm along $T_{uH}z$ is a sum of positive exponentials, hence is
convex and bounded by the larger endpoint norm. The path stays in
the common ball. Pairing the difference of the two central moment
values with $H$ gives
\[
0=2\int_0^1\Omega_t(\xi_H,I\xi_H)(T_{uH}z)\,du.
\]
Taming forces $\xi_H=0$, so the positive transformation fixes the
point. This also gives uniqueness of the positive representative
on the free support.

For coverage, let a small $y=(A,B)\in\cQ$ satisfy
$\rho_t(y)=\tau I$. If $\tau>0$, testing~\eqref{eq:projection-error}
on a vector in $\ker A$ is impossible, so $A$ is injective.
For $\tau<0$, $B$ is surjective. At zero, the same argument gives
$\ker A=\ker B^*$. In each case the orbit has a flat-balanced
representative by Proposition~\ref{prop:finite-geometry}.
For nonzero level the construction preserves the incidence pair
$(\operatorname{im}A,AB)$ or $(\ker B,AB)$, respectively. At zero,
restrict to the common support: two full-support factorizations of
the same $AB$ differ by a general linear transformation on that support.
It stays small: taking traces and using the nuclear norm,
\[
2|\tau|\le(1+\epsilon)|y|^2,\qquad
2\|AB\|_*\le|y|^2.
\]
Its flat squared radius is
$\sum_i\sqrt{\tau^2+4s_i^2}\le(2+\epsilon)|y|^2$.
A fixed shrinking of the actual domain thus gives coverage.

On each fixed-rank piece, choose a smooth support frame and use the
implicit-function theorem. For continuity at rank loss, take a
convergent sequence of flat representatives. The matrices $S$ stay
bounded, and each subsequential limit gives a corrected representative
in the limiting complex orbit. Compact-orbit uniqueness identifies
all such limits. For the inverse, the preceding radius bound makes
the flat representatives bounded; pass to a subsequence and use the
same bounded gauges and compact unitary identifications. Its only
possible limiting orbit is the flat representative of the given
corrected limit. This proves inverse continuity without assuming
continuity of the selected matrices $S$.
\end{proof}

\begin{corollary}\label{cor:quotient-transfer}
Suppose the actual quotient norm on the corrected locus is
\[
Q_{t,y}([w])=\inf_{\eta\in\mathfrak{gl}_2(\CC)}
\mathcal B_{t,y}(w-C_y\eta,w-C_y\eta),
\]
where $\mathcal B$ is a positive Hermitian form on the ambient
coefficient space and
\begin{equation}\label{eq:B-hypothesis}
(1-C(|t|+|y|^2))g^{\rm amb}
\le\mathcal B_{t,y}\le
(1+C(|t|+|y|^2))g^{\rm amb}.
\end{equation}
Then the common-complex-orbit map compares $Q_t$ with
$g_\tau^{\rm q}$ with relative error $O(|t|+|x|^2)$ on every
smooth quotient piece, uniformly as the positive eigenvalues of
$A^*A+BB^*$ tend to zero.
\end{corollary}

\begin{proof}
Differentiate $y=T_{S(x)}x$ on a smooth piece. Its derivative is
$T_Sv$ plus a vector in the full complex vertical tangent. That
vector disappears in the infimum, without any estimate on $dS$.
The linear action $T_S$ carries one complex vertical space onto
the other and distorts squared ambient norms by $e^{\pm2\|S\|}$.
Apply~\eqref{eq:B-hypothesis} before taking the infimum, and use
Theorem~\ref{thm:balancing}. At the flat level the complex-vertical
infimum is the quotient metric, as in Lemma~\ref{lem:outer-model}.
\end{proof}

The additional norm hypothesis in this corollary is essential.
A tame symplectic form alone does not identify the harmonic metric.
We verify that hypothesis for the actual HYM family next.
